让我们从一个简单的问题开始：“什么是初始化（initialization）？”当我们查看 C++ 参考\footnote{\url{https://en.cppreference.com/w/cpp/language/initialization}} 中的定义时，会读到：

\begin{myQuote}
变量在构造时获得其初始值（Initialization of a variable provides its initial value at the time of construction）。
\end{myQuote}

我们可以把这个定义翻译成下面的例子：

\begin{codebox}[cpp]
void foo() {
int x = 42;
// ... 稍后使用 'x' ...
}
\end{codebox}

上面我们写了一个带有局部变量 \texttt{x} 的函数。该变量被声明为 \texttt{int}，并用值 \texttt{42} 进行初始化。这只是你为变量赋予初始值的多种方式之一。下面是更多的选择：

\begin{codebox}[cpp]
struct Point { int x; int y; };
// 声明一个自定义类型
Point createPoint(int x) { return {x, -x}; }
int main() {
int x { 42 };
// 列表初始化
double y = { 100.0 };
// 拷贝列表初始化
auto ptr = std::make_unique<float>(90.5f); // 类型推导 auto
auto z = createPoint(42);
// 通过工厂函数
std::string s (10, 'x');
// 调用构造函数
Point p { 10 };
// 聚合初始化
std::array<float, 100> numbers { 1.1f, 2.2f }; // 数组初始化
// ...
}
\end{codebox}

你还能想出许多其他形式来设置值。我们也可以把这个语法扩展到类的非静态数据成员、\texttt{static} 变量、线程局部变量，甚至动态内存分配上。

理论上，初始化很简单：“把一个值放进新创建变量的内存位置”。然而，这样的操作会牵涉到应用程序代码的许多部分（局部作用域 vs. 非局部作用域 vs. 线程），以及内存中的不同位置（比如栈 vs. 堆）。也正因如此，语法或行为可能会略有不同。

C++ 中至少有如下几种初始化形式：

\begin{itemize}
\item 聚合初始化（aggregate initialization）
\item 常量初始化（constant initialization）
\item 默认初始化（default initialization）
\item 直接初始化（direct initialization）
\item 拷贝初始化（copy initialization）
\item 列表初始化（list initialization）
\item 引用初始化（reference initialization）
\item 值初始化（value initialization）
\item 零初始化（zero initialization）
\item 以及相关的主题，如拷贝消除（copy elision）、静态变量、转换序列、构造函数、赋值、动态内存、存储期等等。
\end{itemize}

虽然这份清单听起来很复杂，但我们会一步步地讲解这些主题，逐步揭示核心概念。随后我们会讨论更高级的例子，看看 C++ 的“内部机制”究竟发生了什么。

虽然我们可以用整数和其他数值类型来解释大部分情况，但最好还是用更贴近实际的东西来讲解。本书从一些基础自定义类型开始，然后讨论它们早期实现中可能遇到的各种问题。之后这些类型会不断扩展，为我们提供更多的背景和令人信服的用例。

\section{从简单类型开始}

在 C++ 中定义类或结构体（struct，即自定义类型）可以让你更自然地建模问题域并解决问题。与其用一堆变量和函数来折腾，不如把它们分组，并提供一致的 API（Application Programming Interface，应用程序编程接口）。C++ 提供了一组内建类型，包括布尔、整型、字符和浮点类型。此外，你还可以使用标准库（Standard Library）中的对象，比如各种容器、\texttt{std::string}、\texttt{std::vector}、\texttt{std::map}、\texttt{std::set} 等等。你可以把这些基础组件收集起来，构建出自己的类型。

为了给我们的主题铺陈背景，我们先从一个表示汽车信息的类型开始，它用于一个汽车列表应用。系统从数据库中读取汽车/卡车的信息，并在应用中显示出来。为了简单起见，这个类型包含四个成员：名称（一个 \texttt{std::string}）、生产年份、座位数量，以及引擎功率。

下面是 \texttt{CarInfo} 类型的第一版代码：

\begin{codeboxt}[cpp]{代码清单 1.1 简单的 CarInfo 结构。在 Compiler Explorer 上运行}
#include <iostream>
#include <string>
struct CarInfo {
std::string name;
unsigned year;
unsigned seats;
double power;
};
int main() {
CarInfo firstCar;
firstCar.name = "Renault Megane";
firstCar.year = 2003;
firstCar.seats = 5;
firstCar.power = 116;
std::cout << "name: " << firstCar.name << '\n';
std::cout << "year: " << firstCar.year << '\n';
std::cout << "seats: " << firstCar.seats << '\n';
std::cout << "power (hp): " << firstCar.power << '\n';
}
\end{codeboxt}

\myGraphic{0.9}{images/fig_p018_1.png}{}

在上面的例子中，我们定义了一个简单的结构体，用于保存 \texttt{CarInfo} 的数据。代码非常简单，存在一些问题，并且遵循 C++03 的风格。在接下来的几章里，我会带你通读这段代码，帮你理解其中的问题以及如何消除它们。我们还会对它进行现代化改造，加入最新的 C++（直到 C++20）特性。

首先：\texttt{name}、\texttt{year}、\texttt{seats} 和 \texttt{power} 被称为非静态数据成员（non-static data members）。\texttt{CarInfo} 类的每个实例都拥有自己的一组这些成员。换句话说，我们把变量分组，以创建问题域中模型的表示。用户自定义类型也可能拥有静态数据成员（static data members），它们是给定类型所有实例之间共享的数据。例如，我们可以设想一个名为 \texttt{numAllCars} 的静态成员变量，用来表示程序中创建的车辆总数。我们会在第 11 章“静态变量”中讨论静态数据成员。

现在，让我们详细研究这段代码。在 \texttt{main()} 函数中，变量 \texttt{firstCar} 的定义与声明：

\begin{codebox}[cpp]
CarInfo firstCar;
\end{codebox}

这被称为默认初始化（default initialization）；由于我们的 \texttt{struct} 很简单，它会让所有内建类型的非静态数据成员都保留为不确定值（indeterminate values）。类似地，对于在函数内声明的简单类型（因为这些变量具有自动存储期（automatic storage duration））²，你也会得到同样（潜在有 bug）的效果：

\footnotetext{与之相反，静态对象和线程局部对象会被零初始化（zero initialized）。}

\begin{codebox}[cpp]
void foo() {
int i;
// 不确定值！
double d;
// 不确定值！
}
\end{codebox}

另一方面，\texttt{std::string} 数据成员 \texttt{name} 会处于空状态（一个空字符串），因为它的默认构造函数（default constructor）会被调用。这一点我们稍后详述。

一旦对象被创建但还未初始化，我们就可以访问它的成员并设置合适的值。默认情况下，\texttt{struct} 的成员是公开访问的（而 \texttt{class} 的成员是私有的）。这样一来，我们可以直接访问并修改它们的值。

\begin{myInfoBox}
\textbf{什么是“自动存储期”？}

程序中的所有对象有四种可能的“存储”方式：自动、静态、线程，或动态。自动意味着内存在作用域开始时被分配，比如在函数中。大多数局部变量都具有自动存储期（automatic storage duration）（除了那些声明为 \texttt{static}、\texttt{extern} 或 \texttt{thread\_local} 的变量）。我们会在关于非局部对象的第 11 章³ 中进一步讨论这一点。
\end{myInfoBox}

\section{置为零值}

在创建了 \texttt{CarInfo} 对象之后，多数数据成员都带有一些不确定值，你可能会感到很不满意。我们可以修正这一点，至少确保数据被置为“零”。请看：

\begin{codeboxt}[cpp]{代码清单 1.2 CarInfo 结构的值初始化。在 Compiler Explorer 上运行}
CarInfo emptyCar{};
std::cout << "name: " << emptyCar.name << '\n';
std::cout << "year: " << emptyCar.year << '\n';
std::cout << "seats: " << emptyCar.seats << '\n';
std::cout << "power (hp): " << emptyCar.power << '\n';
\end{codeboxt}

输出为：

\begin{codebox}[cpp]
name:
year: 0
seats: 0
power (hp): 0
\end{codebox}

用空花括号 \texttt{\{\}} 进行的初始化被称为值初始化（value initialization）；默认情况下（对于内建类型，以及具有既非用户提供的也未被删除的默认构造函数的类），它会把数据置为“零”（会因类型不同而有所调整）。这类似于声明并定义如下变量：


\begin{codebox}[cpp]
int i{};
// i == 0
double d{};
// d == 0.0
std::string s{}; // s 是一个空字符串
int j = {}; // 其它形式的值初始化
std::string str = {}; // ...
\end{codebox}

这一次存储期无关紧要，值初始化对静态、动态、线程局部或自动变量表现一致。对于带有默认构造函数的类型（稍后详述），代码会调用它们；以 \texttt{string s;} 为例，会把它初始化为一个空字符串。

\section{用聚合体进行初始化}

我们的结构体非常简单，对于这类类型，C++ 有专门的规则，可以用所谓的聚合初始化（aggregate initialization）来初始化它们的内部值。这种语法也可以用于数组。下面是一些基本例子：

\begin{codeboxt}[cpp]{代码清单 1.3 聚合初始化的基本语法。在 Compiler Explorer 上运行}
// 数组：
int arr[] { 1, 2, 3, 4 };
float numbers[] = { 0.1f, 1.1f, 2.2f, 3.f, 4.f, 5. };
int nums[10] { 1 }; // 1，其余全为 0
// 结构体：
struct Point { int x; int y; };
struct Line { Point p1; Point p2; };
Line longLone {0, 0, 100, 100};
Line anotherLine = {100}; // 其余置为 0
Line shortLine {{-10, -10}, {10, 10}}; // 嵌套
\end{codeboxt}

总结一下，对于上面的代码：

\begin{itemize}
\item 每个数组元素，或每个非静态类成员，按数组下标 / 在类定义中出现的顺序，从初始化器列表的对应项进行拷贝初始化（copy initialized）。
\item 你可以对数组使用列表初始化；当不提供元素个数时，编译器会推导出元素数量。
\item 如果你在初始化器列表中传入的元素少于数组的元素个数，剩余的元素会被值初始化。对于内建类型，这意味着零值。
\item 对于结构体，你可以使用单个初始化器列表，也可以使用嵌套的列表；展开是递归进行的。
\item 如果你提供的值少于聚合体中的数据成员数量，那么剩余的数据成员（按声明顺序）会被有效地值初始化。
\end{itemize}

第一个要点指出每个元素都是拷贝初始化的。等我们了解了显式构造函数（explicit constructors）之后，会回到这个话题，并解释拷贝初始化与直接初始化在语法上的区别。

对于我们的结构体，我们可以写出如下测试代码：

\begin{codeboxt}[cpp]{代码清单 1.4 CarInfo 结构的聚合初始化。在 Compiler Explorer 上运行}
struct CarInfo {
std::string name;
unsigned year;
unsigned seats;
double power;
};
void printInfo(const CarInfo& c) {
std::cout << c.name << ", "
<< c.year << " year, "
<< c.seats << " seats, "
<< c.power << " hp\n";
}
int main() {
CarInfo firstCar{"Megane", 2003, 5, 116 };
printInfo(firstCar);
CarInfo partial{"unknown"};
printInfo(partial);
CarInfo largeCar{"large car", 1975, 10};
printInfo(largeCar);
}
\end{codeboxt}

输出为：

\begin{codebox}[cpp]
Megane, 2003 year, 5 seats, 116 hp
unknown, 0 year, 0 seats, 0 hp
large car, 1975 year, 10 seats, 0 hp
\end{codebox}

不给出完整的定义，聚合体（aggregate）指的是一个简单类型（或数组），它具有全部公开的数据成员、没有虚函数、也没有用户提供的构造函数。

我们会在本书后续部分讨论聚合体（并看到这类类型的完整定义）。此外，还有专门的一章讲解 C++20 中的聚合体与指定初始化器（designated initializer）。

\section{默认数据成员初始化}

如果你想为自己的数据成员提供某个默认值，该怎么办？通过值初始化，你可以让各种类型得到零值，但有时这可能并不够好。从 C++14 起，我们可以利用非静态数据成员初始化器（Non-static Data Member Initializers，NSDMI，也叫 Default Member Initializers，默认成员初始化器）来为聚合体提供默认值。请看：

\begin{codeboxt}[cpp]{代码清单 1.5 默认成员初始化与聚合体。在 Compiler Explorer 上运行}
#include <iostream>
#include <string>
struct CarInfo {
std::string name { "unknown" };
unsigned year { 1920 };
unsigned seats { 4 };
double power { 100. };
};
void printInfo(const CarInfo& c) { /* */ }
int main() {
CarInfo unknown;
printInfo(unknown);
CarInfo zeroed{};
printInfo(zeroed);
CarInfo partial{"large car", 1975};
printInfo(partial);
}
\end{codeboxt}

打印结果为：

\begin{codebox}[cpp]
unknown, 1920 year, 4 seats, 100 hp
unknown, 1920 year, 4 seats, 100 hp
large car, 1975 year, 4 seats, 100 hp
\end{codebox}

这个语法相当直观：你可以在这个数据成员被声明的地方就初始化它。这能防止你的数据带有不确定值而导致的意外 bug。从例子中可以看到，即便你使用默认初始化或值初始化，数据成员也会得到在 \texttt{struct} 声明中提供的值。如果你在聚合初始化器中给出较少的值，剩余成员会从声明中获取它们的默认值。

严格来说，类内成员初始化器（in-class member initializers）从 C++11 起就已可用，但最初并不支持聚合类型。在本节中，我们只是浅尝辄止地介绍了这项好用的技术。关于这个主题，请参阅专门讲解非静态数据成员初始化的章节。

\section{小结}

在本章中，我们讨论了一些简单的自定义类型，并了解了初始化其数据成员的各种方式。我们从带有不确定值的对象出发，讲到零初始化，然后学习了聚合体，以及提供默认值的技术。

需要记住的要点：

\begin{itemize}
\item 对象和变量的默认初始化（default initialization）会让内建类型得到不确定值，或者默认初始化复杂类型（如 \texttt{std::string} 并将其置为空字符串）。因此，务必保证你的对象和简单变量始终被初始化，这一点至关重要。
\item 对于内建类型，像 \texttt{int x\{\};} 这样的值初始化实际上会对它们产生零初始化的效果，从而让它们为零（以各自的类型而言）。
\item 通过值初始化 \texttt{CarInfo car\{\};}，所有数据成员（对于内建类型）会被零初始化，而对于复杂类型则会被默认初始化。
\item 聚合体是全部公开数据成员的简单类型或数组；我们可以用聚合初始化语法来初始化它们。
\item 得益于类内成员初始化器这一特性，你可以为自己的数据成员提供默认值。
\end{itemize}

\subsection{下一步是什么？}

虽然简单类型很方便，但在 C++ 中，我们常常需要构建大型对象，其数据成员彼此依赖，或者带有不变式（invariants）。在这种情况下，最好把它们隐藏在成员函数之后，并在特定条件下才提供访问。在下一章中，我们会了解类与构造函数。我们也会扩展到目前为止学到的知识。
